% (c) Danny Lee Eldridge · AEGIS Community License · DRAFT v3 (correction) — NOT published; needs Danny's "Publish"
\documentclass[aps,prd,twocolumn,superscriptaddress,nofootinbib,showkeys]{revtex4-2}
\usepackage{amsmath,amssymb,hyperref}
\begin{document}
\title{Flavour at the End of Inflation: Koide's Form, Trimaximal Mixing and a Quantised $\bar\theta$ from the Two Zero-Energy Vacua of an $SL(2,\mathbb Z)$ Plateau}
\author{Danny Eldridge}
\email{dannyeldridge54@gmail.com}
\affiliation{Independent researcher, Oklahoma City, Oklahoma, USA}
\date{19 November 2026 (v3 draft: correction after JUNO)}
\begin{abstract}
In a companion paper~\cite{P1}, inflation is driven by the running of the species scale along the type IIB axio-dilaton $\tau$, with potential $V_0\propto|\nabla\ln E(\tau,\tfrac32)|^2$. That potential vanishes at the $SL(2,\mathbb Z)$ fixed points $\rho=e^{i\pi/3}$ and $i$, and a numerical scan of the fundamental domain finds no other zeros. We ask what the flavour sector looks like if its modular fields come to rest at these two vacua. We report the following results.
(i) Two-field simulations of the end of inflation show that the inflaton itself ends at $\rho$ for every transverse kick tested. The line it falls along has a transverse instability band ($0.96\lesssim y\lesssim1.9$), which makes a ``Y-fork'' towards $i$, but the inflaton crosses it too quickly to be diverted.
(ii) The residual $\mathbb Z_3$ at $\rho$ forces a circulant $\sqrt{M_\ell}$. Its eigenvalues are exactly the Koide--Brannen form $\sqrt{m_k}=a+2|b|\cos(\delta+2\pi k/3)$, and $Q=2/3$ is equivalent to equal norm in the singlet and doublet components. The data give amplitude $\sqrt2$ and $\delta=2/9=Q/3$ to within $1\sigma$. Minimal $A_4$ weight-2 constructions placed exactly at $\rho$ do not produce these numbers.
(iii) Correction (v3): with charged leptons at $\rho$ and neutrinos at $i$, $A_4$ gives TM2, $\sin^2\theta_{12}=0.3408$. JUNO's first data, $0.3092\pm0.0087$~\cite{JUNO}, exclude this at $3.6\sigma$ ($4.6\sigma$ combined with NuFIT). Under modular $S_4$ the \emph{same two vacua} give TM1~\cite{TM1}: the neutrino point $-\tfrac12+\tfrac i2$ is $SL(2,\mathbb Z)$-equivalent to $i$. TM1 predicts $\sin^2\theta_{12}=0.3184$ ($+1.1\sigma$), and the sum-rule branch allowed by NuFIT~6.0 gives $\delta_{CP}=286^\circ$ ($\sin^2\theta_{23}=0.561$) or $262^\circ$ ($0.470$). Of the six images of $i$, three give TM1, three give a pattern with a zero entry (excluded), and none gives TM2.
(iv) At $\rho$ the $A_4$ quark determinant is the real number $27/8$. With CP-symmetric couplings $\bar\theta\in\{0,\pi\}$, but the CKM phase vanishes ($|J|<5\times10^{-15}$). The observed phase therefore needs a quark modulus at a generic point with anomaly-free weights, as in Feruglio, Strumia and Titov~\cite{FST}.
All computations were pre-registered, except the choice of TM1, which was made after JUNO (see the correction note). We state plainly which steps are principles and which are derived.
\end{abstract}
\keywords{flavour; modular symmetry; $A_4$; Koide formula; trimaximal mixing; strong CP; inflation; $SL(2,\mathbb Z)$ fixed points; pre-registration}
\maketitle

\section{Two vacua from one potential}
The species-gradient potential $V_0=V_\infty(\tfrac23)^2y^2|\nabla\ln E(\tau,\tfrac32)|^2$ vanishes wherever $\nabla E=0$. In the fundamental domain this happens at $\rho$ and $i$. Numerically $V_0(\rho)=1.6\times10^{-22}$ and $V_0(i)=2.2\times10^{-22}$, and both points are strict minima, with $V_0\approx8.3\times10^{-8}$ (at $\rho$) and $(2.1$--$42)\times10^{-8}$ (at $i$) at distance $10^{-3}$. The inflaton falls along $x=\tfrac12$, which is a valley near $\rho$ ($y<0.954$) and a ridge for $0.96\lesssim y\lesssim1.9$, where the low side is the line $x=0$ that leads to $i$. We call this the Y-fork. A two-field integration with the hyperbolic kinetic metric shows the inflaton crossing the band and settling at $\rho$ for transverse kicks $\delta x=0$--$3\times10^{-2}$. These kicks are far larger than its quantum spread. Inflation ends at $y\approx2.4$, so the CMB observables of Ref.~\cite{P1} are unchanged. Whether a slow modulus drains to $i$ or to $\rho$ is computed in Sec.~V.

\section{Charged leptons at $\rho$: Koide's form}
At $\rho$ the residual symmetry is $\mathbb Z_3=\langle ST\rangle$. A square-root mass matrix covariant under the cyclic $\mathbb Z_3$ is circulant, $\sqrt M=a\mathbb 1+bC+\bar bC^2$, with eigenvalues
\begin{equation}
\sqrt{m_k}=a+2|b|\cos(\delta+2\pi k/3),\qquad Q=\tfrac13+\tfrac23|b|^2/a^2 .
\end{equation}
Koide's relation $Q=2/3$ is therefore $|b|^2=a^2/2$: the singlet and doublet components of $\sqrt M$ carry equal norm. The algebra is known~\cite{Foot}. What is new here is that inflation selects the $\mathbb Z_3$ point dynamically. Using Belle~II $m_\tau=1777.09\pm0.14$~MeV~\cite{Belle} we find $Q=0.6666735\pm0.0000077$ ($+0.9\sigma$), amplitude $1.414228$, and $\delta=0.2222136\pm0.0000095$ against $2/9$ ($-0.9\sigma$). With $m_\tau$ as the only input, $A=\sqrt2$ and $\delta=2/9$ give $m_\mu=105.666$~MeV and $m_e=0.511029$~MeV. The data also satisfy $\delta=Q/3$ to within $1\sigma$. We do not derive $\sqrt2$ or $2/9$. A locked test of four minimal $A_4$ weight-2 constructions exactly at $\rho$ gives $Q=0.360$, $0.341$, $1/3$ and $1/3$; all fail. The fixed point fixes the \emph{form}, not the numbers.

\section{Neutrinos at $i$: from TM2 to TM1}
At $i$ the residual symmetry is $\mathbb Z_2=\langle S\rangle$. With the $A_4$ triplet of Ref.~\cite{F17}, the $S=+1$ eigenvector fixes the second PMNS column at $(1/3,1/3,1/3)$, which is TM2, with $\sin^2\theta_{12}=1/(3\cos^2\theta_{13})=0.3408$. Version 2 of this note reported this at $+2.8\sigma$ against NuFIT~6.0. JUNO's first measurement, $\sin^2\theta_{12}=0.3092\pm0.0087$~\cite{JUNO}, now places exact TM2 at $+3.6\sigma$ ($+4.6\sigma$ combined), so we withdraw it.

The fixed points are unchanged under modular $S_4\simeq\Gamma_4$~\cite{TM1,NPP4}: charged leptons sit at the $\rho$-image fixed by $ST$, and neutrinos at an image of $i$. The neutrino point used in Ref.~\cite{TM1}, $\tau_{SU}=-\tfrac12+\tfrac i2$, maps to $i$ under $T^{-1}S$, so it is one of the potential's two zero-energy vacua. We computed the fixed column for each of the six conjugates of $S$ in the $S_4$ triplet (24 elements verified). Three give TM1, $(2/3,1/6,1/6)$; three give $(\tfrac12,\tfrac12,0)$, which data exclude; none gives TM2. TM1 gives $\sin^2\theta_{12}=(1-3s_{13}^2)/(3(1-s_{13}^2))=0.3184$, i.e.\ $+1.1\sigma$ from JUNO. Its sum rule $\cos\delta=-\cot2\theta_{23}(1-5s_{13}^2)/(2\sqrt2\,s_{13}\sqrt{1-3s_{13}^2})$ gives $74^\circ$ or $286^\circ$ for $\sin^2\theta_{23}=0.561$, and $98^\circ$ or $262^\circ$ for $0.470$. The first branch lies outside NuFIT~6.0's $3\sigma$ range for both data sets. The prediction is therefore $\delta_{CP}\simeq262^\circ$--$286^\circ$ ($\sin\delta<0$), to be tested by DUNE and Hyper-Kamiokande. The flow is $SL(2,\mathbb Z)$-equivariant, so a slow neutrino modulus that starts uniformly on the level-4 cover reaches each of the six images of $i$ equally often (each conjugate of $S$ arises from 4 of the 24 cosets). Combined with Sec.~V, a fraction $0.69\times\tfrac12=0.345$ of starting points ends in the viable TM1 configuration. Nothing in the potential yet explains why the flavour group has level 4 rather than 3; we state this as an assumption.

\section{Quarks: quantised $\bar\theta$ and the CKM phase}
For the standard $A_4$ structure $\mathrm{diag}(\alpha,\beta,\gamma)\cdot[[Y_1,Y_3,Y_2],[Y_2,Y_1,Y_3],[Y_3,Y_2,Y_1]]$, the determinant at $\rho$ is $Y_1^3+Y_2^3+Y_3^3-3Y_1Y_2Y_3=27/8$, which is real. Both $\rho$ and $i$ are CP-symmetric points. With real couplings $\bar\theta\in\{0,\pi\}$, and $\bar\theta=0$ is a discrete sign choice. In the same setting, however, $|J|<5\times10^{-15}$ over $6000$ samples, against the observed $3.1\times10^{-5}$. The CKM phase therefore requires the mechanism of Ref.~\cite{FST}: anomaly-free quark weights $\sum(2k_Q+k_u+k_d)=0$, so that $\det M_q$ is a real constant, together with a quark modulus at a generic point, for example $\tau_q\approx\tfrac18+i$. In this construction the quark modulus is not one of the two vacua. The use of residual symmetries at fixed points to fix mixing patterns is standard in modular flavour models~\cite{NPP}; what is new here is the dynamical reason for the vacua.

\section{Do slow fields reach $i$?}
A slow modulus does not carry the inflaton's momentum and is overdamped by Hubble friction, so it follows the gradient flow $\dot x=-2y^2\partial_xV$, $\dot y=-2y^2\partial_yV$. We ran this flow under a pre-registered lock (PREREG\_BASINS) from 2000 starting points distributed uniformly in hyperbolic area over the fundamental domain with $y\le3$. A fraction $f_i=0.690$ drains to $i$ and $f_\rho=0.310$ to $\rho$; none ends elsewhere. Along the line $y=3$, starting points with $x\in[0,0.355]$ reach $i$ and those with $x\in[0.358,0.5]$ reach $\rho$. The $\rho$ basin is a strip around the CP-symmetric line $x=\tfrac12$, along which the inflaton falls. A modulus starting at a generic position therefore ends at $i$ about two times in three. The criterion set in advance, $f_i\ge0.25$, is met, so the assignment of the neutrino modulus to $i$ is the natural outcome of the same potential rather than an independent assumption, provided that modulus is slow and obeys the same potential.

\section{What is assumed}
We assume the following. The species-gradient potential itself~\cite{P1}. That the flavour group is modular $S_4$ (level 4), not $A_4$ (this is not derived). Equipartition, $|b|^2=a^2/2$, together with $\delta=Q/3$. That the lepton moduli obey the same potential (the neutrino modulus then reaches $i$ for 69\% of starting points, Sec.~V). The anomaly-free quark weights and the quark modulus at a generic point. None of these is derived here. The vacuum energy beyond zero and the origin of $\Lambda_{\rm obs}$ are not addressed.

\emph{Correction note (v3).} Version 2 predicted exact TM2 and the $\delta_{CP}$ values $236^\circ/125^\circ$ and $287^\circ/74^\circ$. JUNO excludes TM2, and those values are withdrawn. The TM1 replacement was adopted \emph{after} JUNO excluded TM2, so it is post hoc and is not counted as evidence; only its fixed-point counting and its $\delta_{CP}$ scoring were locked before computing (repository D1652, D1667--D1669). The T2K+NOvA joint result~\cite{T2KNOVA} puts $\delta_{CP}$ in $[214^\circ,313^\circ]$ at $1\sigma$, consistent with the $286^\circ$ branch (D1690). The quark sector is hosted on a separate modulus, not predicted (doi:10.5281/zenodo.24399408). \emph{Integrity note.} All computations were pre-registered (locks D1595--D1607, results D1610). Fixes to code were logged before any rerun, and one invalid run is kept on record. The TM1 $\delta_{CP}$ values are statements to be scored against DUNE and Hyper-Kamiokande; because TM1 was chosen after JUNO they will count only as consistency, not as evidence.

\begin{thebibliography}{99}
\bibitem{P1} D.~Eldridge, ``Inflation from the running of the species scale'' (draft, 2026).
\bibitem{FST} F.~Feruglio, A.~Strumia, A.~Titov, JHEP 07 (2023) 027, arXiv:2305.08908.
\bibitem{Foot} R.~Foot, hep-ph/9402242; Y.~Koide, Lett.~Nuovo Cim.~34 (1982) 201; C.~A.~Brannen (2006).
\bibitem{Belle} Belle II, arXiv:2305.19116.
\bibitem{F17} F.~Feruglio, arXiv:1706.08749.
\bibitem{NPP} P.~P.~Novichkov, J.~T.~Penedo, S.~T.~Petcov, JHEP 04 (2021) 206, arXiv:2102.07488; F.~Feruglio, arXiv:2211.00659.
\bibitem{NuFIT} I.~Esteban et al., NuFIT 6.0, arXiv:2410.05380.
\bibitem{JUNO} JUNO Collaboration, first measurement of reactor neutrino oscillations, arXiv:2511.14593.
\bibitem{TM1} Modular TM$_1$ mixing in light of precision measurement in JUNO, arXiv:2511.16348.
\bibitem{NPP4} P.~P.~Novichkov, J.~T.~Penedo, S.~T.~Petcov, A.~V.~Titov, arXiv:1811.04933; G.-J.~Ding et al., arXiv:1910.03460.
\bibitem{T2KNOVA} T2K and NOvA Collaborations, joint oscillation analysis, Nature (2025).
\end{thebibliography}
\end{document}
